\section{Far-region geometry and energy estimates}\label{fe:section}

We prove the energy estimate used in \Cref{sec:far} for a smooth
finite slab $0\le t\le t_{\rm end}$.  The structural inputs are
the lab equation and differentiated gauge in Lemma~\ref{far:structure},
and the initial and starting-annulus bounds in
Lemma~\ref{far:boundary-input}.  We derive the point estimates needed
for the coefficients, then prove the two multiplier estimates conditional
on explicit source budgets.  Appendix~\ref{fs:section} supplies those
budgets; the simultaneous exterior closure is proved in
Proposition~\ref{far:closure}. The radial weighted
currents follow the method of Dafermos--Rodnianski
\cite[Sections~3--4]{DafermosRodnianski2010}, adapted here to the moving emitter
coordinates and the coupled source estimates.

Throughout, smallness of the finite-slab norms is a bootstrap hypothesis;
their finiteness follows from finite-time smooth existence and symbol
decay.  Constants on superlinear terms may depend on all fixed scales.
Linear constants acting \emph{on other far quantities} in absorption or
finite hierarchies are independent of both large matching radii.  Take
$R_0$ beyond all spatial kernels and the transition where target smoothing
reaches its prescribed far form.  Comparable smoothed radii use the same
positive acceleration majorant $\widetilde a$ on a slightly enlarged
sampling interval at that scale.

\subsection{Emitter coordinates and the three energy levels}\label{fe:coordinates}

For the calculations in both far appendices we abbreviate
$T_l,C_e,s_e,\ell_n,\delta_f,\tau_f,l_f,m_f,M_f$ by
$T,C,s,\ell,\delta,\tau,l,m,M$, respectively, and write
$X,\bar X$ for $X_g,\bar X_g$.  These are the same
coordinates, fields and weights as in \Cref{sec:far}.


In the lab chart put
\begin{equation*}
C(q)=-\lambda(q),\quad s(q)=C'(q),\quad t=q+r,\quad z=C(q)+rn,\quad |n|=1,\quad k=1-s(q)\cdot n .
\end{equation*}
These are smooth coordinates for large emitter radius \(r\), by the small speed and the implicit function theorem; the radii \(|x|,|y|\) of the reduction are comparable to \(r\) (the position differences on time intervals of length \(O(r)\) are small times \(r\); use also the contraction defining \(x\)). Work with
\begin{equation*}
T=\partial_{t,{\rm lab}},\quad L=\partial_{r|q,n}=T+n\cdot\partial_z,\quad
 Z=(S=rL,\Omega),\quad s_s=-v_s,\quad D=T+s_s\cdot\partial_z ,
\end{equation*}
where \(\Omega\) are the spherical rotations at fixed \(q,r\). Angular differentiations below can equivalently use finite smooth frames. Components of tensor fields in all the far wave differentiations are \emph{lab} components. Write \(H^{ij}=g^{ij}-\eta^{ij}\), \(\ell_i=(1,-n)\), and subscript \(\ell\ell\) for contraction of upper indices with this pair. The lab null contravariant vector corresponding to \(-\eta^{ij}\ell_j\) is \(L\). Derivatives have the identities
\begin{equation*}
\partial_i=\ell_i T+r^{-1}b_i Z
     =\ell_i k^{-1}\partial_q+r^{-1}(b_i-\ell_i k^{-1}s^d b_d)Z ;
\end{equation*}
\(b_t=0\), \(b_dZ=n_d S+\) a tangential spherical derivative for spatial indices. In particular the \(q\)-derivative of the second good-frame coefficient (before \(r^{-1}\)) is proportional to \(\ell_i\) and of size \(O(|a(q)|)\). All filtered unsmoothed acceleration jets needed below are bounded and decay with at least \(\langle q\rangle^{-A/2}\) on positive times, by the compact budgets and \eqref{co2:sliced-acceleration}, and are zero at negative times. The enlarged average has
\begin{equation*}
r\widetilde a\lesssim \mathsf N(1+t/r)^{-A/2}
\end{equation*}
by \(L^1\) for \(t\lesssim r\) and weighted Cauchy--Schwarz for \(t\gg r\). Here and in low point estimates \(\mathsf N\) can include the data norm, and fixed compact-reaching constants multiplying small quantities are allowed when such bounds are used \emph{nonlinearly}. Derivatives of \(s_s\) of positive order \(d\) are \(O(r^{1-d}\widetilde a)\) in the lab chart; \(B-\eta\) has symbol size \(O(r^{-1})\) there in all the spacetime derivatives. With ordered commutations (\(T^j\) first, then \(Z^I\)) these retain their respective bounds (\(T\) a lab derivative, \(Z\) scaled first order); arbitrary re-expansions differentiating emitter frame factors need not gain powers for additional \(T\)'s. For example \(T^m n=O(1/r)\) for finite \(m\ge1\), with \(Z\)-derivatives allowed.

Use three levels labeled \(0,H,P\). Their unscaled fields \(U_i\), number of commutations \(N_i\), time exponents \(A_i\), and outgoing powers \(p_i\) are
\begin{equation*}
\begin{array}{c|c|c|c|c}
 i&U_i&N_i&A_i&p_i\cr\hline
 0&h&9&0&.64\cr
 H&h&8&1-\eta_*&1-\eta_*/6\cr
 P&\psi=D h&7&A=1.36&1.62 .
 \end{array}
\end{equation*}
Here \(\eta_*>0\) is sufficiently small, say at most \(10^{-8}\). Put \(\delta=\eta_*/100,\ \nu=10\delta,\ e=.1\delta,\ d_*=200\eta_*\).
The strict margins serve different purposes and will be kept throughout:
\begin{center}
\begin{tabular}{@{}ll@{}}
\toprule
Margin & Use \\
\midrule
$p_0,p_H<1$ & Initial energy for arbitrary $r^{-1}$ metric tails \\
$p_i>A_i$, $p_i+.1A_i<2$ & Positive outgoing bulk and Hardy control \\
$A>1$ & Time integrability of the acceleration \\
$2A_H>A$ & Weak-field products in the strong equation \\
$p_P-A=.26$ & Radius-uniform pairing with adjoint zero modes \\
\bottomrule
\end{tabular}
\end{center}
Take
\begin{equation*}
\mu=\langle q/r^\delta\rangle,\quad l=\langle q/r^\nu\rangle,\quad
 m=(r^{-1}+l^{-1})^{-1},\quad W_i=r^{p_i-A_i}(t+10r)^{A_i}.
\end{equation*}
Use time-energy weights, for commutations with \(j\) lab times,
\begin{equation*}
\begin{split}
 w_H&=r^{-e}\mu^{A_H},\\
 w_{P,j}&=w_H l^{A-A_H}m^{-(7-j)d_*},\\
 w_{0,j}&=r^{-.02}\,m_-^{j d_*}m_+^{-(9-j)d_*}.
\end{split}
\end{equation*}
Here \(m_-\) is a smooth positive monotone weight in \(-q/r^\nu\) comparable to \(1+(-q/r^\nu)_+\) (constant on the nonpositive half-line, bounded logarithmic-symbol derivatives); \(m_+\) is the analogous weight in \(q/r^\nu\), capped smoothly by harmonic addition at \(r\) as for \(m\). The comparability in this text always allows fixed constants. Put \(\tau=1+t/r\). Thus \(W_i\asymp r^{p_i}\tau^{A_i}\); \(W_i/w_{i,j}\ll r^{2-.3}\), and in particular \(w_{P,j}\gtrsim \langle t\rangle^A r^{-.025},\ w_H\gtrsim \langle t\rangle^{A_H} r^{-e-\delta A_H}\) for \(t\gg r\).

Our scalar wave unknowns (componentwise) are
\begin{equation*}
\chi=Z^I(rT^j U_i),\qquad n_*=|I|+j\le N_i .
\end{equation*}
We write \(B_1\chi\) for its full first derivatives (size comparable to \(|\partial_q\chi|+r^{-1}|Z\chi|\)), and \(G_1\chi=r^{-1}Z\chi\). In formulas indicating bounds, lists of fields/derivatives denote their Euclidean size, with sums over finite components or lower commutations as specified. The main measured quantities, before specifying a current correction, are
\begin{equation}\label{fe:F1}
\begin{array}{ll}
 L^2(dt\,dr\,d\omega/r): & W_i^{1/2}(G_1\chi,\chi/r),\quad
                       w_{i,j}^{1/2}B_1\chi ,\\
 L^\infty_t L^2(dr\,d\omega):& W_i^{1/2} G_1\chi,\quad
                       w_{i,j}^{1/2}B_1\chi
 \end{array}
\end{equation}
on \(r\ge R_0\). Here \(d\omega\) is round measure of \(n\). In particular the \emph{flat} good derivatives are included. The total size \(\mathsf F\) sums these and the extra current norms below; quantities on a starting overlap annulus a fixed factor farther in, when needed in estimates, can instead use \(C(R_0)\mathsf K_R+C_{\rm fix}(e_{\rm in}+\mathsf N^2)\), by the co/lab comparisons on bounded annuli. Indeed changing \(u\) to \(\psi\) costs only field-derivative products with strong coefficient differences such as \(v(t)-v_s\) there. No supremum of characteristic flux in this overlap is required from compact estimates. All a priori integrations are for smooth solutions with finite symbol decay on bounded intervals; or one first truncates at infinity where this decay gives the limiting estimates. The time-zero energies have the stated orders by the initial symbol controls (\(T^j\) costs at least \(r^{-j}\) on undifferentiated initial jets, startup flat). For \(h\), even with no lab-time commutation, \(p_i<1\); for \(D h\) the extra decay easily gives finiteness. Constants on the input may depend on fixed scales.

\subsection{The wave-gauge constraint and weak-null structure}\label{fe:weak-null}

The weak-null organization of the Einstein equations in wave coordinates
is due to Lindblad--Rodnianski \cite[Sections~3--5]{LindbladRodnianski2003}; see also
\cite{LindbladRodnianski2010}.
The component calculation below records the precise polarization and
weighted source structure required here.


The lab gauge identity on the end gives an improved differentiated constraint. With \(\Gamma_\mu(g)=g^{ab}\Gamma(g)_{\mu;ab}\) (semicolon here separates the lowered first Christoffel index), it is
\begin{equation*}
\Gamma_\mu(g)=g_{\mu\sigma}g^{ab}\Gamma(B)^\sigma_{ab}+\Theta_\mu .
\end{equation*}
In particular its size is \(O(r^{-2})\). Subtracting the first derivatives of \(B\), the flat linear gauge of \(h\) costs terms with smooth coefficients of sizes
\begin{equation}\label{fe:F2}
O((r^{-1}+|h|)\partial h)+O(r^{-2}h)+\Theta .
\end{equation}
These are structural product bounds: the first coefficient is smooth vanishing in \(B-\eta,h\); linear-size fixed coefficients obey lab symbol differentiation, including the enlarged-average bounds for \(\Theta\); differentiation in \(h\) obeys the chain rule. Also the reduced wave equation has the form
\begin{equation*}
g^{ij}\partial_{ij}h=f_h ;
\end{equation*}
its pure-background term is of symbol size \(O(r^{-2}\widetilde a)\), with \(D\) derivative \(O(r^{-3}\widetilde a)\), and the remaining terms have sizes/structure
\begin{equation}\label{fe:F3}
O(r^{-2}\partial h+r^{-3}h)+Q_\eta(\partial h,\partial h)
                    +O((r^{-1}+|h|)|\partial h|^2).
\end{equation}
Here smooth bounded multiplicative coefficients, lab symbol background coefficients as above, and smooth functions of \(h\) are understood in the \(O\)'s. \(Q_\eta\) is a constant flat quadratic polynomial. On replacing \(\partial_i h\) by \(\ell_i p\) with \(\ell^\mu(p_{\mu\nu}-\frac12\eta_{\mu\nu}{\rm tr}_\eta p)=0\), it is a constant multiple of
\begin{equation*}
\ell\otimes\ell\ |\Pi(n)p|^2 ,
\end{equation*}
where \(\Pi\) takes the two-dimensional tracefree tangential part on the spatial screen. Thus errors of this replacement with \(p=T h\) use either a good derivative or a factor of the gauge error \eqref{fe:F2}. To check the structure, write Ricci as \(-\tfrac12 g^{ab}\partial_{ab}g_{\mu\nu}+\partial_{(\nu}\Gamma_{\mu)}(g)\) plus the products
\begin{equation*}
(\partial_a g^{ab})\Gamma(g)_{b;\mu\nu}
 -(\partial_{(\nu}g^{ab})\partial_a g_{\mu)b}
 +\Gamma^a_{a d}\Gamma^d_{\mu\nu}-\Gamma^a_{\nu d}\Gamma^d_{a\mu}.
\end{equation*}
Substitute the gauge identity. Pure terms at \(h=0\) are Ricci of the target plus differentiated compensator terms. In the flat product evaluation the first two displayed products cancel on \(\ell p\) as above. Writing the Christoffel matrices at fixed \(\mu\) as \((\ell_\mu p^\sharp+p_\mu^\sharp\otimes\ell-\ell^\sharp\otimes p_\mu)/2\), the last two products then give \((({\rm tr}\,p)^2/8-{ \rm tr}(p^2)/4)\ell_\mu\ell_\nu\), as asserted by the constrained-component relation in the flat null frame. This calculation also works polarized. One can differentiate in lab coordinates first and project afterwards, or retain \eqref{fe:F2} as an exact error identity; when differentiating projected formulas each lab differentiation hitting the frame has a factor \(O(1/r)\) with bounded emitter jets. Here the tangent \(p^\sharp\) notation uses flat raising.
For completeness, rotate the screen so that $n=(1,0,0)$ and use
$\eta=\diag(-1,1,1,1)$, $\ell=(1,-1,0,0)$.  The four gauge relations
are solved by the six-parameter matrix
\begin{equation*}
 p=\begin{pmatrix}
 a&b&c&d\\
 b&-a-2b&-c&-d\\
 c&-c&e&f\\
 d&-d&f&-e
 \end{pmatrix}.
\end{equation*}
Consequently $\tr_\eta p=-2(a+b)$,
$\tr((\eta^{-1}p)^2)=2(a+b)^2+2(e^2+f^2)$, and
\begin{equation*}
 \frac18(\tr_\eta p)^2-\frac14\tr((\eta^{-1}p)^2)
       =-\frac12(e^2+f^2)=-\frac14|\Pi p|^2.
\end{equation*}
Polarization gives the same identity with $|\Pi p|^2$ replaced by the
screen pairing of two constrained tensors.  Thus the quadratic restriction
and the projection cancellation remain valid when the two factors carry
different ordered commutations.

We will also use the corresponding projection onto the constrained plane.
Define
\begin{equation}\label{fe:gauge-projection}
\begin{aligned}
 \mathcal G_np&=\ell^\mu
       (p_{\mu\nu}-\eta_{\mu\nu}\tr_\eta p/2),
 &\breve p&=p-J_n\mathcal G_np,\\
 \mathcal G_nJ_n&=\Id,
 &\Pi J_n&=0.
\end{aligned}
\end{equation}
The matrix calculation shows that $\mathcal G_n$ has rank four and that
every screen polarization occurs in its kernel.  Its restriction to
$\ker\Pi$ is therefore surjective; taking the Euclidean right inverse
on this restriction gives a smooth bounded $J_n$.  Consequently
$\mathcal G_n\breve p=0$ and $\Pi\breve p=\Pi p$.  The four gauge
errors measure exactly the $LL$, $L$--screen, and screen-trace components.


\subsection{The scalar emitter operator}\label{fe:emitter-operator}

For a scalar, the emitter operator is
\begin{equation*}
P_g\chi:=kr g^{ij}\partial_{ij}(\chi/r),\qquad
 P_\eta\chi=-2\partial_q\partial_r\chi+r^{-2}\mathcal A_s\chi ,
\end{equation*}
\begin{equation}\label{fe:F4}
\mathcal A_s=(1+s\cdot n)(S^2-S)-2(s\cdot n)(S-1)
       +2s^A\nabla_A(S-1)+ k\Delta_{\mathbb S^2}.
\end{equation}
Here is a direct computation that also records the absence of acceleration
terms in the flat operator.  Put $\sigma=s\cdot n$ and let $\gamma$ be the
round screen metric.  The nonzero covariant entries are
\begin{equation*}
 \eta_{qq}=-1+|s|^2,\quad \eta_{qr}=-k,\quad
 \eta_{qA}=rs_A,\quad \eta_{AB}=r^2\gamma_{AB}.
\end{equation*}
Their inverse and density are
\begin{equation*}
 \eta^{qr}=-k^{-1},\quad \eta^{rr}=(1+\sigma)/k,\quad
 \eta^{rA}=s^A/(kr),\quad \eta^{AB}=r^{-2}\gamma^{AB},\quad
 \sqrt{|\det\eta|}=kr^2\sqrt{\det\gamma},
\end{equation*}
with $\eta^{qq}=\eta^{qA}=0$.
In the divergence formula the only coefficient differentiated in $q$ is
$kr^2\eta^{qr}=-r^2$, which is independent of $q$.
Using $\nabla_As^A=-2\sigma$ and $\nabla_Ak=-s_A$ gives
\begin{equation*}
 P_\eta\chi=-2\chi_{qr}+(1+\sigma)\chi_{rr}
 -\frac{2\sigma}{r}\chi_r+\frac{2\sigma}{r^2}\chi
 +\frac{2s^A}{r}\nabla_A\chi_r
 -\frac{2s^A}{r^2}\nabla_A\chi+\frac{k}{r^2}\Delta_\gamma\chi.
\end{equation*}
Replacing $r\partial_r$ by $S$ proves \eqref{fe:F4}. More generally for an upper-index coefficient \(J^{ij}\), the additional expression \(P[J]\chi=kr J^{ij}\partial_{ij}(\chi/r)\) consists of
\begin{equation}\label{fe:F5}
J_{\ell\ell}\partial_q(k^{-1}\partial_q\chi)
  + O\left(
    \frac{J}{r}(1,Z)\partial_q\chi+
    \frac{J}{r^2}(1,Z,Z^2)\chi+
    \frac{a(q)}{r}J_{\ell\ell}(1,Z)\chi\right),
\end{equation}
where the other coefficients are smooth bounded functions of \(n,s(q)\). This follows from the derivative formula above (including differentiating \(1/r\)); a derivative of a good-frame coefficient in \(q\) has the extra \(\ell\) contraction. The term with \(a(q)\) allows contractions of its components. These identities apply \emph{after} lab times have been commuted.

\subsection{Spacelike anchors and point estimates}\label{fe:anchors}

The following estimates are implications of the finite-time bootstrap norms, \eqref{fe:F2}, and \eqref{fe:F3}; they do not assume the uniform estimate that will close the bootstrap.

The starting scale can and will always be taken so large that even the inner band radii of the following auxiliary construction (a fixed constant times \(Q\gtrsim R_0^{.45}\)) are in the lab symbol regime beyond all spatial kernels and all fixed smoothing transitions. No comparison with compact estimates at radii of this size is used as input for the spacelike observation.

\begin{lemma}[Spacelike anchors]\label{fe:spacelike-anchors}
For \(r\asymp R_f,\ q\le-R_f^\theta/10\), \(\theta=.45\), the center is zero. On a band \(-q\asymp Q\), we have
\begin{equation}\label{fe:F6}
\|Z^I(rT^j h)\|_{L^2_\omega}\lesssim (e_{\rm in}+\mathsf S_a)(r/Q)^{1/2}
               \quad (|I|+j\le9),
\end{equation}
allowing fixed input constants; the constant here on \(\mathsf S_a\) is uniform far out. Also the ordinary unweighted \(L^2_z\) energy on such a band at fixed \(t\) of Minkowski commutations through nine is \(\lesssim (e_{\rm in}+\mathsf S_a)Q^{-1/2}\).
\end{lemma}

\begin{proof}
These anchors will be proved without a characteristic commutator estimate.
Up to time \(\lesssim R_f\) work outside
\begin{equation*}
|z|\ge t+cQ+C_1\int_0^t ((u+Q)Q)^{-1/2}du
\end{equation*}
with \(c\) a sufficiently small fixed band constant, \(C_1\) large fixed. Indeed its ratio to $Q$ is bounded by
$C\sqrt{R_f}Q^{-3/2}\lesssim R_f^{(1-3\theta)/2}=R_f^{-.175}$,
so the correction is $o(Q)$ uniformly on the band. Bootstrap unweighted slice energies of \(h\) after rotations, boosts, translations and the flat dilation through order nine by a large constant times \((e_{\rm in}+\mathsf S_a) Q^{-1/2}\). The initial norm involves at most ten derivatives of $h$ and at most
nine derivatives of its initial normal jet, all controlled by the small
initial seminorm.  The squared integral of the leading
$r^{-2}$ first jet outside $r\asymp Q$ is $O(Q^{-1})$, which gives the
claimed initial energy. Radial integration to infinity gives the commuted fields in \(L^2_\omega\) of size \(O(|z|^{-1/2})\) times the energy, and sphere sup with two spare rotations. The boundary moves at least as fast outwards as the outgoing light speeds including bootstrap errors. The lab time multiplier for the true principal is uniformly timelike and controls full unweighted first derivatives. Here are the energy error details. Commuting the principal leaves, besides constant linear multiples of lower flat wave terms (eliminate recursively), terms \(\mathcal Z^b H^{ij}\partial^2\mathcal Z^d h\), \(b+d\le9,d<9\), for the Minkowski list \(\mathcal Z\), including index mixing by commutators with derivatives. The first factor is bounded on the sphere by \(C/|z|+C_{\rm boot}(e_{\rm in}+\mathsf S_a)(Q|z|)^{-1/2}\); if high, use only \(L^2_\omega\) on it, placing two extra rotations on the low second factor. Smooth composition has the same placements. One derivative of the second factor costs at most \(C/Q\) times one dilation/boost/rotation commutation and lower terms (apply to first derivatives). Derivative-of-principal energy errors and the differentiated quadratic first-jet terms of \eqref{fe:F3} obey the same product bounds, using the derivative placement with \(1/Q\) on a coefficient. Linear lower terms with symbol-decaying coefficients do also, by energy/Hardy or the commuted field bound. For instance a coefficient \(O(|z|^{-3})\) times a high field costs \(C(t+Q)^{-2}\) times the energy at that order. The time integrals of all multiplier losses just described are \(o(1)\), including \(O(\sqrt{R_f}Q^{-3/2})\). The pure defect commutations cost \(C\mathsf S_a |z|^{-3}\) pointwise and \(\lesssim\mathsf S_a Q^{-1/2}\) in \(L^1_t L^2_z\). This proves the auxiliary bootstrap by the ordinary domain energy inequality. At a spacelike evaluation band, \(Z\) and its iterates can be expanded with bounded coefficients in these commutations (\(r(T+\partial_{|z|})\) uses the sum of a boost and dilation); the corresponding first-jet statements include \(O(1/r)\) coefficient derivatives. This gives \eqref{fe:F6}, uniformly also if the lifespan truncates the band.
\end{proof}

In low point estimates on products below we abbreviate a constant times small bootstrap size by \(N\). This can include fixed input and fixed compact-to-far point constants. When \eqref{fe:F6} is used on a \emph{linear} energy error we keep its specified independent bounds instead. Endowing point bounds with \(L^2_\omega\) means they hold on each sphere.

\begin{lemma}[Radial point estimates]\label{fe:point-estimates}
The finite-time norms \eqref{fe:F1} imply
\begin{equation}\label{fe:F7}
\begin{array}{ll}
 \|\chi(h)\|_{L^2_\omega}\lesssim N r^{2\eta_*}\tau^{-A_H/2}
                          & (n_*\le8),\\
 \|\chi(h)\|_{L^2_\omega}\lesssim N r^{.52}
                          & (n_*\le9),\\
 \|B_1\chi(h)\|_{L^2_\omega}\lesssim N r^{e/2}\mu^{-A_H/2}
                          & (n_*\le7).
 \end{array}
\end{equation}
\end{lemma}

\begin{proof}
For \(t\lesssim r\) integrate to an anchor with \(r'\asymp r\) large enough compared with \(t\), using \eqref{fe:F6}; the integral of \(w_H^{-1}\) costs \(\lesssim r^{e+\delta A_H+1-A_H}\). For \(t\gg r\) integrate inward to the fixed starting annulus, using a compact weighted endpoint anchor there (an average suffices). The top time weight is everywhere at least a constant times \(r^{-.021}\). For the last estimate take a radial unit trace using the endpoint bad derivatives also at the next order, and (if undifferentiated terms occur) the first bound divided by \(r\). In detail, \(Z,\partial_q\) commute and \(\partial_q(rT^j h)=k r T^{j+1}h+r^{-1}s^d b_dZ(rT^j h)+O(r^{-1})rT^j h\). Thus \(\partial_q^2\) in this trace expands into first derivatives of the next and lower fields with bounded coefficients and possible \(r^{-1}\) times fields through order eight. The harmless field terms satisfy the required weights, and first derivatives against \(w_H\) have a comparable weight throughout the interval. Angular and \(r|q\) derivatives are simpler. This proves \eqref{fe:F7}; the compact overlap or one-sided traces suffice near a radial cutoff. The same convention of two spare rotations gives angular sup versions. In particular for \(n_*\le7\),
\begin{equation}\label{fe:F8}
\|\chi(\psi)\|_{L^2_\omega}
           \lesssim N r^{e/2}\mu^{-A_H/2}.
\end{equation}
Use \(D=(1-s_s\cdot n)T+r^{-1}s_s^d b_dZ\), with lab-time commutators of \(s_s\) taken \emph{before} expansion, and \eqref{fe:F7}; \(rZ^IT^j\) or \(Z^I rT^j\) field lists are equivalent. All first-derivative versions at one lower order follow too, up to harmless \(O(1/r)\) lower terms.
\end{proof}

\subsection{Improved null projections}\label{fe:projections}

\begin{lemma}[Differentiated null projection]\label{fe:null-projection}
In any bound with a \(T^j\) on an upper-index tensor, contract with \(\ell\ell\) \emph{after} those derivatives and then apply the indicated \(Z\)'s to the full contraction.
\begin{equation}\label{fe:F9}
\begin{array}{ll}
 \|Z^I H_{\ell\ell}\|_{L^2_\omega}
    \lesssim r^{-1}+N r^{-1+2\eta_*}\tau^{-A_H/2}&(|I|\le8),\\
 \|Z^I (T^j H)_{\ell\ell}\|_{L^2_\omega}
    \lesssim r^{-2}+N r^{-2+8\eta_*}\tau^{-A_H/2}
                                 &(j\ge1,\ |I|+j\le8),\\
 \|\partial_q Z^I (T^j H)_{\ell\ell}\|_{L^2_\omega}
    \lesssim r^{-2}+N r^{-2+8\eta_*}\tau^{-A_H/2}
                                 &(|I|+j\le7).
 \end{array}
\end{equation}
\end{lemma}

\begin{proof}
The first estimate follows by smooth inversion of $B+h$, the symbol
bound on $B-\eta$, and \eqref{fe:F7} through order eight.
For the last two, \eqref{fe:F2} after ordered differentiation controls the bad derivative of $h_{LL}$, the $L$--screen components, and the spherical trace of $h$: \eqref{fe:F7} bounds each good derivative by
\begin{equation*}
 C N r^{-2+2\eta_*}\tau^{-A_H/2},
\end{equation*} and products such as \(H\,\partial h\) satisfy the bound as well (lab differentiations may be distributed before projecting; the highest differentiated \(\partial h\) is of total order at most eight and its bad part can use a \(T\)-field). This is just solving the flat gauge on the \(\ell\) gradient using \eqref{fe:F2} and \eqref{fe:gauge-projection}, using \(r^{-1}b Z\) for the omitted part, or using \(k^{-1}\partial_q\) for the bad derivative. Derivatives of bounded metric compositions have at most one high sphere factor; inverse contraction differs from the negated flat-raised covariant perturbation by a smooth quadratic. \(\Theta\) costs at most \(C\widetilde a/r\) with ordered derivatives allowed. The \(r^{-2}\) term at this precision is for the stationary background contribution, not a claimed \(\tau\)-decay of all of \(H\).
\end{proof}

\begin{lemma}[Tight cone projection and top-order control]\label{fe:tight-projection}
Set $b=1+\min(|q|,r)$.  Then
\begin{equation}\label{fe:F10}
\begin{aligned}
    r\|Z^I H_{\ell\ell}\|_{L^2_\omega} &\lesssim 1
                   &&(|I|\le8,\ |q|\le 10 r^\theta),\\
    r\|Z^I (T^j H)_{\ell\ell}\|_{L^2_\omega} &\lesssim r^{-.46}
                   &&(j\ge1,\ |I|+j\le9),\\
    r\|Z^9 H_{\ell\ell}\|_{L^2_\omega}
       &\lesssim \log r + b^{1/2}r^{-p_0/2}+b r^{-1+8\eta_*}.
 \end{aligned}
\end{equation}
Constants here can be bounded independently of \(R_0\) by choosing smallness of the nonlinear quantities after the radii; equivalently retain background plus \(O(C_{\rm fix}\mathsf N)\) in them.
\end{lemma}

\begin{proof}
The second follows exactly as in \eqref{fe:F9}, using the rough \(r^{.52}\) order-nine field bound. For the others it suffices to prove corresponding covariant \(h_{LL}\) estimates: smooth quadratic inverse errors at order nine after multiplication by \(r\) are even \(O(r^{-.46})\).

\smallskip\noindent\emph{The constrained radial integral.}
Keep \(t,n\) fixed and integrate the equation for the bad derivative of \(Z^d(r h_{LL})\) from negative \(q\) toward the cone. At fixed \(t\) this adds a flat good derivative, since the integration differentiates by \(\partial_q-\partial_r\). By \eqref{fe:F2} the terms to be integrated have the following sphere sizes, after distributing \(Z\)'s through order \(d\): flat \(G_1\chi(h),\chi(h)/r\); products of \(B_1\chi(h),\chi(h)/r\) with factors from the order \(r^{-1}\) background perturbation or \(r^{-1}\chi(h)\), with further bounded metric compositions allowed; and \(O(\widetilde a)\). Here all \(\chi\)'s need only spatial-emitter commutations (no lab times); this follows using \(\partial_i=(\ell_i/k)\partial_q+\) good and \(Z\) on the full constrained combinations. 

\emph{Negative-$q$ anchor bands.}
When \(t\asymp r\), start at \(-q\asymp r\); radii during integration stay comparable up to a target in \(t\asymp r\). In the spacelike bands \(-q\asymp Q\gtrsim r^\theta\), for \(d\le8\) the integral of good terms costs \(CN (Q/r)^{1/2}\) bandwise by \eqref{fe:F6}. Products there are summable as well: a bad derivative in these products through that order is bounded using the extra Minkowski commutation with loss \(C/Q\), giving e.g. \(O(N/Q)\) for the quadratic band cost. 
For \(d=9\), use the unweighted spacelike-band slice energy for a highest first derivative, giving \(O(N)\) per band; undifferentiated order-nine fields use \eqref{fe:F6}, and in products with a high field the bad derivative now costs \(1/Q\) only on a low factor. 

\emph{Near the cone.}
Up through \(|q|\lesssim r^\theta\) the remaining integral of \(G_1\chi\) costs \(\lesssim N r^{(\theta-p_i)/2}\) at the corresponding \(H,0\) time-slice orders. For \(d=9\), \(\chi_9/r\) with nine \(Z\)'s uses the \(W_H\)-controlled \(G_1\chi_8\) on this slice. Products with high \(B_1\chi_9\) use \(L^2_{q,\omega}\) size \(\lesssim N r^{.011}\) on comparable radii and gain the factor \(r^{-1+2\eta_*}\); products with \(\chi_9/r\) and low \(B_1\chi\) can use the same \(G_1\chi_8\) bound and \eqref{fe:F7} on the low factor. All lower undifferentiated terms alone cost at most \(CN r^{-1+2\eta_*}\) per unit length. These estimates give tightness for \(d\le8\), at most logarithmic loss for nine. 

\emph{Longer intervals and timelike evaluation.}
For evaluation up to length \(b\gtrsim r^\theta\) they cost at most the two additional terms of the last bound, using time-slice Cauchy--Schwarz with length \(b\). Targets with \(t\ll r\) need only \eqref{fe:F6}. If \(t\gg r\), start instead by a length-\(\asymp r\) local average of \(Z^9(rh_{LL})\) on comparable radii (use the \(G_1\chi_8\) time-slice norm), then integrate the same estimates over this length. This proves \eqref{fe:F10}.
\end{proof}

\subsection{A uniform bound for the lowest radiative jets}\label{fe:news}

\begin{lemma}[Uniform news bound]
On every finite smooth bootstrap interval, the Cartesian lab-time derivatives satisfy
\begin{equation}\label{fe:F11}
 |r\Pi T h|+|r\Pi T^2h|+|r\Pi T\psi|\lesssim N.
\end{equation}
The constant uses only the previously stated low bootstrap, initial, and compact overlap controls.
\end{lemma}

\begin{proof}
Put $\chi_j=rT^jh$, $j=0,1$, and
$\mathcal Y_j=\Pi k^{-1}\partial_q\chi_j$.
The projection $\Pi$ depends on $n$ only; $k$ depends on $q,n$ only.
Thus projection of the two leading terms of the scalar emitter operator is exactly
\begin{equation}\label{fe:news-principal}
 -2\Pi\partial_q\partial_r\chi_j
 +H_{\ell\ell}\Pi\partial_q(k^{-1}\partial_q\chi_j)
 =(-2k\partial_r+H_{\ell\ell}\partial_q)\mathcal Y_j.
\end{equation}
We show that every remaining term is bounded by $CNr^{-2+.04}$.
This calculation uses only low commutations, so it precedes and supplies the
undifferentiated supremum used in the high-order products of \Cref{fs:section}.

Choose a temporary fixed $\varepsilon>0$ larger than the finitely many
$\eta_*,e$ losses in the following bounds, with $3\varepsilon<.04$.
The sphere estimates in \eqref{fe:F7}, with two spare rotations, give
\begin{equation*}
 |Z^{\le2}\chi_j|+|Z^{\le1}\partial_q\chi_j|
       +|\partial_q^2\chi_0|\lesssim Nr^{\varepsilon},
       \qquad j=0,1.
\end{equation*}
For the last term, use
$\partial_q(rT^jh)=krT^{j+1}h+s^db_dZ(T^jh)$,
differentiate once, and commute $\partial_q$ through $Z$.
This leaves first derivatives of the next and lower fields,
$r^{-1}$ times good derivatives, and $a(q)\partial_q\chi_0$.
The largest field count in the displayed point estimates, including the two
rotations, is at most $j+2+2\le5$; the unit trace for a bad derivative uses
one further order, still below the weak cap $N_H=8$.
Also
\begin{equation*}
 |H|\lesssim r^{-1}+Nr^{-1+\varepsilon},\qquad
 |H_{\ell\ell}|\lesssim r^{-1+\varepsilon},\qquad |a(q)|\lesssim N.
\end{equation*}
The angular term $r^{-2}\Pi\mathcal A_s\chi_j$ is therefore
$O(Nr^{-2+\varepsilon})$.
The three remainder types in \eqref{fe:F5} have bounds
\begin{equation*}
 \frac{|H|}{r}|(1,Z)\partial_q\chi_j|,\quad
 \frac{|H|}{r^2}|(1,Z,Z^2)\chi_j|,\quad
 \frac{|a(q)H_{\ell\ell}|}{r}|(1,Z)\chi_j|
       \lesssim Nr^{-2+2\varepsilon}.
\end{equation*}
In particular no uncontrolled radial power accompanies the acceleration.
One also has
\begin{equation*}
 r|a(q)H_{\ell\ell}|\lesssim1,
\end{equation*}
by \eqref{fe:F10} near the cone and the decay of $a(q)$ away from it.

For $j=1$, commuting the Cartesian derivative $T$ through the physical wave
operator introduces $-P[TH]\chi_0$.  The full coefficient obeys
$TH=O(r^{-2})+O(Nr^{-1+\varepsilon})$, whereas \eqref{fe:F9} gives
\begin{equation*}
 (TH)_{\ell\ell}=O(r^{-2})+O(Nr^{-2+\varepsilon}).
\end{equation*}
Thus its two-bad part, $(TH)_{\ell\ell}\partial_q(k^{-1}\partial_q\chi_0)$,
is $O(Nr^{-2+2\varepsilon})$.
The mixed part has the additional $r^{-1}$ in \eqref{fe:F5}; its good and
acceleration parts are smaller.  It is the improved double contraction,
rather than an improved bound on every component of $TH$, that is needed.

For the quadratic source, polarize the constant Cartesian polynomial before
projecting.  In particular
$TQ_\eta(\partial h,\partial h)=2Q_\eta(\partial Th,\partial h)$.
For $h_i=T^ih$, $i=0,1$, write
\begin{equation*}
 \partial_\alpha h_i=\ell_\alpha\breve p_i+E_{\alpha i},
 \qquad
 \ell^\mu\bigl((\breve p_i)_{\mu\nu}
       -\tfrac12\eta_{\mu\nu}\tr_\eta\breve p_i\bigr)=0,
 \qquad \Pi\breve p_i=\Pi T^{i+1}h.
\end{equation*}
Take $\breve p_i$ from \eqref{fe:gauge-projection} with
$p=T^{i+1}h$.  The good derivative terms and
\eqref{fe:F2}, differentiated once in Cartesian lab time when $i=1$, give
$|E_i|\lesssim Nr^{-2+2\varepsilon}$.
For example the differentiated gauge error contains $(Th)\partial h$,
of size $O(N^2r^{-2+2\varepsilon})$, and not a term of order $r^{-1}$.
The constrained--constrained polarization is a multiple of
$\ell\otimes\ell$ and is annihilated by $\Pi$.
Every other quadratic term has one $E_i$ and one full first jet of size
$O(Nr^{-1+\varepsilon})$; after the physical-source factor $r$, it is
$O(N^2r^{-2+3\varepsilon})$.
Differentiating an already expanded frame instead gives $T\ell=O(r^{-1})$
and the same suppression.  The vanishing-coefficient quadratic terms in
\eqref{fe:F3}, including their single $T$ derivative, also satisfy this bound.
The linear lower terms are $O(Nr^{-2+\varepsilon})$, and the pure defect
costs $r^{-1}\widetilde a\lesssim Nr^{-2}$.
We have proved
\begin{equation}\label{fe:news-transport}
 |(-2k\partial_r+H_{\ell\ell}\partial_q)\mathcal Y_j|
       \lesssim Nr^{-2+.04}.
\end{equation}

At fixed $n$, this transport direction has strictly decreasing lab time, with
$dr/d\sigma=-2k<-c$ and
$dt/dr=1-H_{\ell\ell}/(2k)=1+O(r^{-1+2\eta_*})$.
Its backward integral curve reaches $t=0$ or the fixed-radius starting region
and stays within the current finite lifespan.
Since $\int_{R_0}^{\infty}r^{-2+.04}\,dr<\infty$, the initial symbol bounds
and compact low point bounds control $\mathcal Y_j$ by $CN$.
The displayed identity for $\partial_q\chi_j$ converts this to the first
two terms of \eqref{fe:F11}, the additional term being good.
Finally differentiate
$\psi=(1-s_s\cdot n)Th+r^{-1}s_s^db_dZh$ once in lab time, taking the
lab derivative before expanding the frame.  It gives
$r\Pi T\psi=(1-s_s\cdot n)r\Pi T^2h$ plus good-frame and
$\widetilde a$ terms bounded by the same low estimates.
This proves the last term of \eqref{fe:F11}.
\end{proof}

\subsection{True null directions, spacelike caps, and corrected currents}\label{fe:currents}

We now construct the energy currents and specify the source norms that they require. Let \(X,\bar X\) in this calculation denote the \emph{true} outward/inward future null normals to the emitter spheres, normalized by \(dt=1\), and \(G=-g(X,\bar X)>0\). These are algebraic smooth functions of \(g,n\) near flat, with components in the lab chart, independent of \(s(q)\) given those arguments. They are \(O(r^{-1}+|h|)\) away from the corresponding flat normals. With \(z_X=n\cdot X^z\), one has \(Xq=(1-z_X)/k\), bounded in size by a constant times
\begin{equation*}
d_g=|H_{\ell\ell}|+|H|^2.
\end{equation*}
Indeed the linearization of the radial speed on the outgoing root costs just the \(L L\) covariant variation. Also \(Xr=1-Xq\). The estimates above imply \(d_g\lesssim r^{-1}\) for \(|q|\le 2r^\theta\), \(d_g\lesssim r^{-1+2\eta_*}\) elsewhere. Thus we can use global spacelike caps \(\mathfrak u=u\), where for example
\begin{equation*}
\mathfrak u=q+\int_{R_0/2}^r
       [\rho^{\delta-1}+\rho^{20\eta_*-1}j(|q|/\rho^\theta)]\,d\rho ,
\end{equation*}
with \(j=1\) above 2 and zero below 1, a smooth transition. One has \(\partial_q\mathfrak u=1+o(1)\), \(X\mathfrak u>0,\ \bar X\mathfrak u>c_0>0\) for \(R_0\) large, and \(|\mathfrak u-q|\lesssim r^\delta+|q|^{20\eta_*/\theta}\). The induced changes between \(r^\delta+|q|\) and \(r^\delta+|u|\) are comparable. No angular component of \(d\mathfrak u\) on the spheres occurs. Fix
\begin{equation*}
M=r^{.72}\mu^{-A_H}j_0(q/r),\qquad K=1+r^{-1}\int_q^\infty M(q',r)\,dq',
\end{equation*}
with \(j_0\) nonnegative, one on \([-1,1]\), zero outside \([-2,2]\), smooth. Then \(K=1+o(1)\), \(r|XK|=o(1)\), \(-r\bar X K\gtrsim M-o(1)\). Indeed the integral costs \(o(r)\) with the same bounds after \(r\partial_{r|q}\), and \(M Xq=o(1)\).

For each equation a specified scalar current correction \(\mathcal P\) will be used (possibly zero), with
\begin{equation*}
\partial_q=\alpha\bar X+\text{a combination of }X,\ e_A,\quad
  c=\frac{G\alpha}{2k},\qquad
  \vartheta=-g(\bar X,\cdot)/G,\qquad
  \xi=d\chi+c\mathcal P\vartheta .
\end{equation*}
Here \(e_A\) is a \(g\)-orthonormal screen, the indicated coefficients are bounded, and \(\alpha=(z_X-s\cdot n)/(z_X-z_{\bar X})\). Complete the definition of \(\mathsf F\) by including
\begin{equation}\label{fe:F12}
\|(W_i(1+M))^{1/2}\xi_X\|_{L^2(dt\,dr\,d\omega/r)}
       +\sup_u\|W_i^{1/2}\xi_X\|_{L^2(dr\,d\omega;\ \mathfrak u=u)}
\end{equation}
out to the current time. Only the directional component, not \(M^{1/2}\) times the uncorrected good frame, is required with the extra \(M\).

Use the notation \(\|\cdot\|_{\rm b}\) here just for \(L^2(dt\,dr\,d\omega/r)\). The source splitting in the outgoing multiplier may be written
\begin{equation}\label{fe:F13}
P_g\chi=\lambda_1 L\chi/r+ f_{\rm flat}
                     +\partial_q\mathcal P+f_{\rm r}+f_{\rm mix},
\end{equation}
where \(\lambda_1\ge0\) is bounded and the flat triangular terms are specified below. In the lab-time multiplier we keep the whole raw source \(f_{\rm raw}\) (after the first two terms). The budgets to estimate for each level are
\begin{equation}\label{fe:F14}
\begin{split}
 &\|W_i^{1/2}(\mathcal P,Z\mathcal P,r d_g\,\partial_q\mathcal P)\|_{\rm b}
       +\sup_t\|W_i^{1/2}\mathcal P\|_{L^2_{r,\omega}},\\
 &\| (W_i/(1+M))^{1/2}r f_{\rm r}\|_{\rm b}
       +\|w_{i,j}^{1/2}r f_{\rm raw}\|_{\rm b}.
 \end{split}
\end{equation}
For \(f_{\rm mix}\), supported where \(t\lesssim r\), an alternative charge on each dyadic shell \(r\asymp R_f\) is as follows. Suppose
\begin{equation*}
\|r f_{\rm mix}\|_{L^2_\omega}
       \le l_{R_f}(\mathfrak u)\,\mathcal T(t,r),
            \qquad \sup_t\|\mathcal T(t,\cdot)\|_{L^2_r}
                          \le E_{\rm infield}.
\end{equation*}
The corresponding mixed budget is
\begin{equation}\label{fs:mixed-charge}
 \sum_{R_f}R_f^{(p_i-1)/2}
       \|l_{R_f}\|_{L^2(du)}E_{\rm infield}.
\end{equation}
To see the shell power, pair the source with $W_i\xi_X$ using measure
comparable to $dt\,dr\,d\omega$.  Cauchy--Schwarz on each outgoing cap
and then in $u$ gives
\begin{equation*}
 C R_f^{p_i/2-1}
 \sup_u\|W_i^{1/2}\xi_X\|_{L^2(\mathfrak u=u)}
 \|l_{R_f}\|_{L^2(du)}
 \left(\int_{t\lesssim R_f}\!\!\int\mathcal T(t,r)^2\,dr\,dt\right)^{1/2}.
\end{equation*}
The last factor is at most $C R_f^{1/2}E_{\rm infield}$, producing
exactly the stated charge. Positive sums of such charges are allowed.

\begin{proposition}[Far energy reduction]\label{fe:energy-reduction}
Assume that each equation has the splitting \eqref{fe:F13}, and that its correction, raw, remainder, and mixed budgets satisfy, in the triangular order specified below, bounds of the form
\begin{equation}\label{fe:source-hypothesis}
o_{R_0\to\infty}(1)(\mathsf F+\mathsf S_a)
  +C\sum_{\rm prior}\mathsf F_{\rm prior}
  +C(R_0)\mathsf K_R+C_{\rm fix}(e_{\rm in}+\mathsf N^{1+c_1})
\end{equation}
for fixed \(c_1>0\). Assume also that the corrections have initial weighted energy and averaged
starting-annulus costs bounded by the initial and overlap terms in
\eqref{fe:source-hypothesis}; the annular norm is the compact weighted
$L^2$ norm. Then the two energy multipliers give
\begin{equation}\label{fe:far-conclusion}
\mathsf F\ \le\ o(1)(\mathsf F+\mathsf S_a)
          +C(R_0)\mathsf K_R+C_{\rm fix}(e_{\rm in}+\mathsf N^{1+c_2})
\end{equation}
(with \(c_2>0\), allowing analogous small input bootstrap errors). It suffices also to take tunably small uniform coefficients on \(\mathsf F\) in this argument. A convenient prior order for linear terms is: level \(P\) before the \(h\) hierarchies, within a level lower total orders first, more lab times first at fixed total order, and then fewer \(S\)'s first. Small radius gains and perturbative nonlinear effects may cross this order. Norms aggregated at equal indices are treated together.
\end{proposition}

\begin{proof}
\par\medskip\noindent\textbf{Claim 1: the corrected divergence identity.}\quad
Use volume \(r^{-2}d{\rm vol}_g\), comparable to \(dt\,dr\,d\omega\); prime on a divergence indicates this density. The normalized equation for \({\rm div}'(d\chi)^\sharp\) differs from \(P_g\chi/k\) by a first-order coefficient \(O(r^{-2})\) (contracted Christoffels) and a potential of size
\begin{equation*}
O(r^{-2})\,O\big(|s(q)|+|H|+r|a(q)H_{\ell\ell}|\big).
\end{equation*}
The latter follows from \eqref{fe:F4}, \eqref{fe:F5} on a constant and is small times \(r^{-2}\), by \eqref{fe:F9}, \eqref{fe:F10} and the acceleration decay. Write \(Q(\xi)=\xi\otimes\xi-g^{-1}(\xi,\xi)g/2\). For a multiplier \(Y\) the primed divergence identity has deformation part
\begin{equation*}
g(\nabla_{\xi^\sharp}Y,\xi^\sharp)
                          -\tfrac12{\rm div}'Y\,g^{-1}(\xi,\xi)
\end{equation*}
and other terms \(\xi_Y\,{\rm div}'\xi^\sharp+d\xi(\xi^\sharp,Y)\).
Use $Y=W_iKX$ and factor out its scalar weight.  To calculate the
correction, write $F=c\mathcal P$.  Since
$\vartheta^\sharp=-\bar X/G$, we have
\begin{equation*}
 \xi^\sharp=-\frac{\xi_{\bar X}}G X-\frac{\xi_X}G\bar X
                 +\xi_{\rm scr}^{\sharp},\qquad
 \vartheta(X)=1,\quad \vartheta(\xi^\sharp)=-\xi_{\bar X}/G.
\end{equation*}
The differentiated-$F$ contribution from divergence is
$-G^{-1}\xi_X\bar XF$.  In the curl contribution the term
$-G^{-1}\xi_{\bar X}XF$ from $dF(\xi^\sharp)$ cancels
$-dF(X)\vartheta(\xi^\sharp)$ exactly.
The sum is therefore
\begin{equation*}
 -2G^{-1}\xi_X\bar XF
       +\nabla_{\rm scr}F\cdot\xi_{\rm scr}.
\end{equation*}
Since $2c/G=\alpha/k$, this cancels the indicated source component
with its exact sign.  When the derivative falls on $\vartheta$, the divergence term is
proportional to $\xi_X$, while in
$F\,d\vartheta(\xi^\sharp,X)$ the $\xi_{\bar X}$ part of
$\xi^\sharp$ vanishes by $d\vartheta(X,X)=0$.
Thus no remaining correction term from divergence or curl involves
$\xi_{\bar X}$. They are bounded, without the weight, by
\begin{equation*}
C r^{-1}
  \big[
  |\xi_X|((1+\sqrt M)|\mathcal P|+|Z\mathcal P|+r d_g|\partial_q\mathcal P|)
     +|\xi_{\rm scr}|(|\mathcal P|+|Z\mathcal P|)\big].
\end{equation*}
To check coefficients, first derivatives in a good true direction of \(g,n,X,\bar X\) in lab components cost \(O(1/r)\); first unrestricted derivatives cost \(O((1+\sqrt M)/r)\), by \eqref{fe:F7} and the frame formulas. The same estimates apply to \(\vartheta\), in particular to \(d\vartheta(e_A,X)\) and \(d\vartheta(\bar X,X)\), respectively. Screen derivatives of \(c\) cost \(O(1/r)\); for \(\bar X c\) the extra dependence on \(s(q)\) uses \(\alpha/k=(1+(z_X-1)/k)/(z_X-z_{\bar X})\), hence \(r|a(q)|d_g\lesssim1\). Also \(\partial_q-\alpha\bar X\) differentiates \(\mathcal P\) by true good directions, giving the \(r d_g\) allowance.

The outgoing deformation has a lower bound of a positive constant times
\begin{equation}\label{fe:F16}
\frac{W_i}{r}\big((1+M)|\xi_X|^2+|\xi_{\rm scr}|^2\big)
\end{equation}
up to small errors in the joint estimates. For details, the derivative of \(W_iK\) alone gives
\begin{equation*}
 -G^{-1}\bar X(W_iK)|\xi_X|^2-\tfrac12 X(W_iK)|\xi_{\rm scr}|^2.
\end{equation*}
Here \(-r\bar XW_i/W_i\) is positive with fixed margin (each \(p_i>A_i\)), and \(r XW_i/W_i<2\) with fixed margin. The following frame estimates check the remaining deformation:
\begin{equation*}
\begin{split}
  &g(X,\nabla_Y X)=0,\\
  &g(e_B,\nabla_{e_A}X)=\delta_{AB}/r+O(r^{-2+.04}),\\
  &g(e_B,\nabla_X X),\ g(\bar X,\nabla_X X)=O(r^{-2+.04}),\\
  &{\rm tr}_{\rm scr}\nabla X-2Xr/r=O(r^{-2+.04}).
 \end{split}
\end{equation*}
Indeed differentiation of the algebraic normals in \(X\) gains \(Xg=O(r^{-2+.04})\), \(Xn=O(r^{-2+.04})\); spherical differentiation gives the ordinary flat expansion, and Christoffels with the displayed slots use good derivatives or the improved first derivative of \(g_{LL}\) from \eqref{fe:F9}, to the precision indicated. The other components of \(\nabla X\), after subtracting the flat emitter ones of order \(1/r\), cost \(O(r^{-2+.04})+O(N\sqrt{1+M}/r)\). Any non-small flat off-expansion components relevant to cross terms between \(\xi_X,\xi_{\rm scr}\) or to \(\xi_X^2\) here cost only \(O(|s|/r)\) (for the flat \(L\), incoming variation of \(n\) is \(O(|s|/r)\)). Consequently the potentially two-bad square vanishes exactly by nullity, mixed terms involving the bad derivative have coefficients \(O(r^{-2+.04})\), and the screen expansion supplies the missing positive \(W_iK|\xi_{\rm scr}|^2/r\). For example the square of a mixed coefficient, relative to the two
positive bulk weights, is bounded by
$r^{-2+.08}W_i/w_{i,j}\lesssim r^{-.22}$.
We absorb mixed bad terms using \(W_i/w_{i,j}\ll r^{1.7}\); other cross terms use smallness and \eqref{fe:F16}. Here \(\xi_{\bar X}=\bar X\chi\).

The favorable volume sign is positive with the future energy convention (time endpoint energy is subtracted at positive divergence). Integrate on slabs cut off either just by a time endpoint or also by an outgoing cap. All future endpoint stresses in this argument are nonnegative as energies. At a time endpoint the outgoing one controls \(W_i\) times \(|\xi_X|^2+|\xi_{\rm scr}|^2\); at \(\mathfrak u=u\) it controls \(W_i|\xi_X|^2\,dr\,d\omega\).
To verify the latter with its measure, $d\mathfrak u$ has no screen
component, so a future normal is
\begin{equation*}
 -\nabla\mathfrak u=\frac{\bar X\mathfrak u}{G}X
                  +\frac{X\mathfrak u}{G}\bar X.
\end{equation*}
The flux density after the coarea Jacobian is, up to the uniformly
comparable primed volume factor,
\begin{equation*}
 \frac{W_iK}{\partial_q\mathfrak u}
 \left(\frac{\bar X\mathfrak u}{G}|\xi_X|^2
                +\frac{X\mathfrak u}{2}|\xi_{\rm scr}|^2\right)
       \,dr\,d\omega.
\end{equation*}
Here $\partial_q\mathfrak u\asymp1$, $\bar X\mathfrak u\ge c_0$ and
$X\mathfrak u>0$.  This proves the uniform lower bound required in
\eqref{fe:F12}, even when $X\mathfrak u$ tends to zero at infinity.
 The starting radial flux costs are estimated by averaging the lower radius in \([R_0/2,R_0]\) and using the stated compact comparisons. The Hardy control of \(\chi/r\) in \eqref{fe:F1} follows within this argument: \({\rm div}' X=O(r^{-2+.04})\), so integrate \({\rm div}'(W_i X|\chi|^2/r^2)\), using the \emph{negative} strict leading derivative of \(W_i/r^2\); only past/radial anchors are charged and \(X\chi=\xi_X-c\mathcal P\).
In detail, $X\log(W_i/r^2)\le-\epsilon_i/r$ for some fixed
$\epsilon_i>0$, and hence
\begin{equation*}
 \operatorname{div}'(W_iX|\chi|^2/r^2)
 \le -\frac{\epsilon_i}{2}\frac{W_i|\chi|^2}{r^3}
            +\frac{2W_i}{r^2}\chi X\chi.
\end{equation*}
Young's inequality bounds the last term by half the displayed negative
bulk plus $C_iW_i|X\chi|^2/r$.  Future boundary terms of this scalar
current are nonnegative with the energy convention, so the bulk
$W_i^{1/2}\chi/r$ norm is bounded by the good bulk, the correction
budget, and the past and radial anchors.
 Conversion from true corrected good derivatives to the flat \(G_1\) in bulk and at time endpoints similarly charges \eqref{fe:F14} and a small amount of bad energy (\(d_g^2 W_i\ll w_{i,j}\)); it is not needed on outgoing caps.
The pointwise comparison behind this step is
\begin{equation*}
 |G_1\chi|\le C(|\xi_X|+|\xi_{\rm scr}|+|\mathcal P|)
                      +C d_g|B_1\chi|.
\end{equation*}
The screen bases span the same emitter sphere.  The coefficient of the
incoming derivative in the difference between $L$ and $X$ is $O(d_g)$;
all other difference coefficients are $O(|H|)$.
Since $d_g^2W_i/w_{i,j}\lesssim r^{-.3+4\eta_*}$, the last term is
small in both the spacetime and the time-endpoint norms.  The correction
terms at time endpoints use the supremum in \eqref{fe:F14}.
 The first-order and potential differences from \(P_g/k\) listed above are absorbable by these estimates. Also \(\xi_X\lambda_1 L\chi/(kr)\) has good sign up to such conversions (its \(\mathcal P\) part costs \eqref{fe:F14}, not claimed small by sign). The remaining wave source terms use exactly \eqref{fe:F14}, \eqref{fe:F16} and the mixed entry.

\par\medskip\noindent\textbf{Claim 2: the lab-time multiplier controls the bad derivatives.}\quad
For the lab-time multiplier \(wT\), \(w=w_{i,j}\), use the uncorrected \(Q(d\chi)\). Its endpoint energy controls the full bad-weight cap on time slices. In the deformation, the two-bad coefficient from \(\mathcal L_T g\) is improved by \eqref{fe:F9}; bad-good mixed errors from metric/density differentiation are absorbable since \(w/W_i\) has a strict negative radial power and \(r|\partial g|\lesssim r^{-1}+N r^{e/2}\mu^{-A_H/2}\) at the required first order. More explicitly \(w_H/W_H\lesssim r^{-e-(p_H-A_H)}\), and the gaps at the other two levels are larger. The weight-derivative part is
\begin{equation*}
-G^{-1} (Xw)Q(d\chi)(\bar X,T)
                     -G^{-1}(\bar Xw)Q(d\chi)(X,T).
\end{equation*}
We have \(r Xw/w<-e/2\): differentiate explicitly, using the tight \(Xq=O(1/r)\) near \(|q|\lesssim r^\theta\) and the looser bound elsewhere. Possible adverse log derivatives from \(m^{-d}\) or \(m_+^{-d}\) in the radial powers are at most \(9d_*\nu\); top \(m_-\) is favorable at fixed \(q<0\). Thus we gain the two-bad bulk with a fixed margin. Adverse \(\bar Xw\) at top is at most \(Cw/r\) by \(q\)-monotonicity; at strong level it is at most \(Cw((r^\delta+|q|)^{-1}+r^{-1})\ll W_P/r\). At weak level the same absorption into \(W_H/r\) holds except \(0\le q\le r^{.3}\); indeed \(1-p_H<.3(1-A_H)\). Here and in this multiplier conversions charging \eqref{fe:F14} are understood.

On the exceptional weak-weight range use \(G_1\chi=r^{-1}Z\chi\), where the undivided \(Z\chi\) is a field of total order at most nine. At \(r\asymp R_f\), uniformly from \(q=-R_f^\theta\) to \(R_f^{.3}\) at fixed \(t\asymp R_f\), \eqref{fe:F6} and the top \(w_{0,j}\)-endpoint give the squared sphere bound
\begin{equation*}
\|Z\chi\|^2_{L^2_\omega}\lesssim
 (e_{\rm in}+\mathsf S_a)^2 R_f^{1-\theta}
       +\mathsf F^2 R_f^{\theta+.021}.
\end{equation*}
On the shell the bad exceptional integral for the flat good derivatives
therefore costs at most
\begin{equation*}
 C(e_{\rm in}+\mathsf S_a+\mathsf F)^2
 R_f^{-1+.3+\max(1-\theta,\theta+.021)}.
\end{equation*} The true-minus-flat derivative contribution is absorbed directly by \(w_H\) bulk with the extra \(d_g^2\). The exponent is explicitly
$-1+.3+\max(.55,.471)=-.15$.
Thus these shell costs are summable with radius gain, explaining why no
extra positive-$q$ weighted bad flux is required at the weak level. All raw wave sources in this multiplier now use \eqref{fe:F14}; the signed leading term in \eqref{fe:F13} just uses the flat good/bad weight gap here.

\par\medskip\noindent\textbf{Claim 3: the flat commutators are triangular.}\quad
To specify the flat terms, commute \(Z\) after lab differentiations, eliminating \(P_g\) at lower orders; separate the \(P[H]\) commutations, including any scalar commutator multiples of it, into the nonflat sources below. The flat recurrence is \(P_\eta S=(S+1)P_\eta+r^{-2}\mathcal A_s\); rotation commutator errors are \(r^{-2}O(s)(1,Z,Z^2)\). Thus each \(S\) count contributes \(+(1+s\cdot n)L\chi/r\) from the pure \(S^2\) leading term on the right, giving \(\lambda_1\) in \eqref{fe:F13}. Terms where extra rotations differentiate its \(s\)-dependent coefficient can instead be counted as \(O(s)\). All remaining size-one terms in \(f_{\rm flat}\) cost, after multiplication by \(r\), good derivatives/field-over-radius of lower orders or of the same order with fewer \(S\)'s, with unchanged \(j\); feedback terms at this precision have \(O(s)\) coefficients. In particular a pure \(r^{-2}\chi\) term at the same order from a commuted \(S\) uses \(r^{-1}\chi=L\chi_{\rm prior}\). This proves their triangular acceptability.

\par\medskip\noindent\textbf{Claim 4: the boundary costs use the stated inputs.}\quad
In this argument entries involving the outer cap supremum need only be estimated beyond \(R_0\). In the averaging annulus of the initial radial cutoff all field bulk/time-slice controls can instead use the compact estimates with constants depending on \(R_0\), which have full mixed orders there. For instance every common principal commutator puts order at most \(N_i+1\) on the differentiated field (one lower commutation count before the two wave derivatives); in the \(D\)-equation an additional order-two term on \(h\) uses at most nine derivatives but carries either a parameter acceleration coefficient or a strong field derivative. Thus field source/current products there can use ordinary time-weighted Sobolev estimates off trapping, not an outgoing cap input on the overlap. On linear pure-defect forcing involving acceleration alone one keeps the direct far \(\widetilde a\) bound with its same radius gain on \(\mathsf S_a\), \emph{not} a large overlap coefficient times \(\mathsf S_a\). The required direct source estimate is proved in \Cref{fs:section}. The surface cutoffs themselves can be selected by averaging against the full time annular bounds, irrespective of the outer cap used. Divergence integrations at infinity can also first take a sequence of radial truncations at each bounded time range; ordinary high symbol decay there suffices on the raw jets (with \(p_H<1\), and an extra decay power for the strong field), and the displayed cap/bulk bounds on the specific corrections, or their bounded-time symbol estimates, suffice for correction fluxes.

\par\medskip\noindent\textbf{Claim 5: finite hierarchical absorption.}\quad
Add the two multiplier estimates with fixed coercive relative factors,
use Hardy and the true-to-flat conversions, and take separate suprema
over the endpoints.  Enumerate the finitely many blocks in the stated
prior order as $1,\ldots,J$ and let $E_k$ denote their full norms.
After Young's inequality each block has an estimate of the form
\begin{equation*}
 E_k^2\le \epsilon(R_0)\sum_{l=1}^JE_l^2
       +C\sum_{l<k}E_l^2+B^2,
\end{equation*}
where $B$ is the nonhierarchical right-hand budget in
\eqref{fe:source-hypothesis}; the same argument applies to tunably
small uniform errors.  All size-one constants $C$ and the number $J$
are independent of the two large matching radii.
Choose fixed positive weights $a_k$ decreasing so rapidly that
$C\sum_{k>l}a_k\le a_l/4$ for every $l$.
Summing $a_kE_k^2$ absorbs every prior-order term.
Only after these weights have been chosen do we enlarge $R_0$ so that
$\epsilon(R_0)\sum_k a_k\le\frac14\min_k a_k$.
The resulting weighted norm is equivalent to the finite sum defining
$\mathsf F$, with constants independent of the matching radii, and
its square root gives \eqref{fe:far-conclusion}. Small multiples of any same-time right hand budgets in the bootstrap can be kept on the right before combining all levels. Compact boundary/source squares costing \(C(R_0)\mathsf K_R^2\) are harmless with these conventions. The source hypotheses are verified in \Cref{fs:section}, yielding \eqref{fs:F}. This proves the stated energy reduction.
\end{proof}
